\section{Method}
Let $\theta_0$ denote a base model and $e_t$ an edit producing $\theta_t$ from $\theta_{t-1}$. The step delta is $\Delta\theta_t=\theta_t-\theta_{t-1}$, and $y_t\in\mathbb{R}^K$ is its vector of benchmark outcomes. We propose predicting $\Delta y_t=y_t-y_{t-1}$ from weights and predicting $y_t$ along edit sequences. For deployed bf16 weights $W$ and proposed update $\Delta$, decision D1 defines $\Delta_{\mathrm{eff}}=\operatorname{bf16}(W+\Delta)-W$. Features should describe this effective change; rounding statistics record lost entries and the norms before and after rounding. An fp32 merge is a separate experiment.

The map $\phi_{\mathrm{stats}}$ collects per-tensor Frobenius and relative norms, cosine similarity of previous and new weights, mean, standard deviation, maximum absolute value, positive and nonzero fractions, leading singular values, spectral entropy and stable rank. The spectral summaries use the retained singular values. The map $\phi_{\mathrm{svd}}$ expresses each matrix update as singular triplets. For LoRA factors it uses QR decompositions followed by an SVD of the small core, removing dependence on the choice of factorization; dense updates use a truncated SVD. The sign of each pair is fixed by the largest absolute component of the left vector. Each rank token contains its singular value and fixed random projections of its left and right vectors. Degenerate singular subspaces may still admit multiple bases.

We propose six predictor families. M1 is a stateless tensor-token Transformer; M2 is a state-space sequence model taking pooled delta features and edit descriptors; M3 predicts a subsequent weight embedding with a JEPA objective and a variance regularizer. M4 is an LLM reader given textual statistics and singular components. M5 is a causal trajectory Transformer, and M6 uses tabular JEPA pretraining on statistics followed by a frozen encoder and supervised head. Delta-norm linear regression, a statistics-based LightGBM model and an SVD-feature MLP provide baselines. These are planned comparisons, not reported predictor results.

We will evaluate family-held-out, seed-held-out and ancestry-grouped splits. Ancestry grouping uses the full manifest so that intermediate unevaluated edits do not separate descendants from their root. Separate family and seed splits do not by themselves guarantee ancestry isolation. We will report per-benchmark Spearman correlation and mean absolute error, with bootstrap confidence intervals over edits; trajectory comparisons will resample independent ancestry groups. The preregistered sequence test asks for at least 15\% relative MAE reduction on trajectories of length at least four.
